% ============================================================
% Prezentacja: Zakres badań – praca magisterska
% Temat: Wpływ strategii zastępowania pikseli na metryki
%        Insertion/Deletion w wyjaśnialności CNN
% ============================================================
\documentclass[aspectratio=169, 12pt]{beamer}

% ── Pakiety ───────────────────────────────────────────────
\usepackage[T1]{fontenc}
\usepackage[utf8]{inputenc}
\usepackage[polish]{babel}
\usepackage{lmodern}
\usepackage{amsmath, amssymb}
\usepackage{graphicx}
\usepackage{booktabs}
\usepackage{multirow}
\usepackage{xcolor}
\usepackage{tikz}
\usepackage{pgfplots}
\pgfplotsset{compat=1.18}
\usepackage{array}
\usepackage{colortbl}
\usepackage{algorithm2e}

% ── Motyw ─────────────────────────────────────────────────
\usetheme{Madrid}
\usecolortheme{default}

\definecolor{mainblue}{RGB}{25, 60, 120}
\definecolor{accentred}{RGB}{180, 30, 30}
\definecolor{lightgray}{RGB}{240, 240, 245}
\definecolor{darkgray}{RGB}{80, 80, 80}
\definecolor{successgreen}{RGB}{30, 130, 60}

\setbeamercolor{frametitle}{bg=mainblue, fg=white}
\setbeamercolor{title}{bg=mainblue, fg=white}
\setbeamercolor{block title}{bg=mainblue!85, fg=white}
\setbeamercolor{block body}{bg=lightgray}
\setbeamercolor{section in toc}{fg=mainblue}
\setbeamercolor{itemize item}{fg=mainblue}
\setbeamercolor{itemize subitem}{fg=accentred}
\setbeamercolor{enumerate item}{fg=mainblue}

\setbeamertemplate{navigation symbols}{}
\setbeamertemplate{footline}{%
  \leavevmode%
  \hbox{%
    \begin{beamercolorbox}[wd=.4\paperwidth,ht=2.5ex,dp=1.125ex,leftskip=.3cm]{author in head/foot}%
      \usebeamerfont{author in head/foot}\insertshortauthor
    \end{beamercolorbox}%
    \begin{beamercolorbox}[wd=.4\paperwidth,ht=2.5ex,dp=1.125ex,center]{title in head/foot}%
      \usebeamerfont{title in head/foot}\insertshorttitle
    \end{beamercolorbox}%
    \begin{beamercolorbox}[wd=.2\paperwidth,ht=2.5ex,dp=1.125ex,rightskip=.3cm,right]{date in head/foot}%
      \insertframenumber{} / \inserttotalframenumber
    \end{beamercolorbox}%
  }%
}

% ── Metadane ──────────────────────────────────────────────
\title[XAI – Insertion/Deletion]{Wpływ strategii zastępowania pikseli\\
na metryki Insertion/Deletion\\
w wyjaśnialności CNN}
\subtitle{Zakres badań eksperymentalnych\\[4pt]
\small Określenie metod i kryteriów porównawczych}
\author[J. Kala]{Jacek Kala}
\institute{Studia magisterskie Data Science}
\date{Listopad 2026}

% ============================================================
\begin{document}
% ============================================================

% ── Slajd tytułowy ────────────────────────────────────────
\begin{frame}
  \titlepage
\end{frame}

% ── Spis treści ───────────────────────────────────────────
\begin{frame}{Plan prezentacji}
  \tableofcontents
\end{frame}

% ============================================================
\section{Motywacja i luka badawcza}
% ============================================================

\begin{frame}{Kontekst: wyjaśnialność modeli CNN}
  \begin{columns}[T]
    \begin{column}{.55\textwidth}
      \begin{itemize}
        \item Modele CNN osiągają wysoką dokładność klasyfikacji obrazów,\\
              ale są \textbf{czarnymi skrzynkami} (\textit{black boxes})
        \item Metody XAI (\textit{Explainable AI}) tworzą mapy istotności
              wskazujące, \textbf{które piksele} wpłynęły na decyzję modelu
        \item Kluczowe pytanie: jak \emph{wiernie} odzwierciedlają
              rzeczywistą logikę modelu?
      \end{itemize}
      \vspace{6pt}
      \begin{block}{Metryki ewaluacji XAI}
        \textbf{Insertion} i \textbf{Deletion} (Petsiuk et al., 2018) \\
        to dominujące metryki jakości map istotności
      \end{block}
    \end{column}
    \begin{column}{.42\textwidth}
      \begin{tikzpicture}[scale=0.85]
        \draw[fill=mainblue!15, rounded corners=4pt] (0,0) rectangle (4,2.4);
        \node[font=\small\bfseries, mainblue] at (2,2.1) {Obraz wejściowy};
        \draw[->, thick, mainblue] (2,0) -- (2,-0.8);
        \draw[fill=mainblue!30, rounded corners=4pt] (0.2,-0.8) rectangle (3.8,-2.0);
        \node[font=\small] at (2,-1.4) {CNN (ResNet50)};
        \draw[->, thick, mainblue] (2,-2.0) -- (2,-2.8);
        \node[font=\small, accentred] at (2,-3.1) {Klasa + mapa istotności};
        \draw[fill=accentred!10, rounded corners=4pt] (0,-3.5) rectangle (4,-4.5);
        \node[font=\footnotesize] at (2,-4.0) {Metoda XAI\\(GradCAM, LIME, \ldots)};
      \end{tikzpicture}
    \end{column}
  \end{columns}
\end{frame}

\begin{frame}{Metryki Insertion i Deletion — definicja}
  \begin{columns}[T]
    \begin{column}{.48\textwidth}
      \begin{block}{Deletion}
        Zaczynamy od \textbf{oryginalnego} obrazu.\\
        Stopniowo \textbf{usuwamy} piksele (sortowane malejąco
        według ważności).\\[4pt]
        \textcolor{accentred}{$\downarrow$ konfidencja = dobra mapa}\\[4pt]
        \[ \text{AUC}_\text{del} = \int_0^1 f_\theta(\hat{x}(t))\,dt \]
        \small im \textbf{mniejsze} AUC, tym lepsza mapa istotności
      \end{block}
    \end{column}
    \begin{column}{.48\textwidth}
      \begin{block}{Insertion}
        Zaczynamy od \textbf{zastąpionego} obrazu (baseline).\\
        Stopniowo \textbf{wstawiamy} piksele (sortowane malejąco
        według ważności).\\[4pt]
        \textcolor{successgreen}{$\uparrow$ konfidencja = dobra mapa}\\[4pt]
        \[ \text{AUC}_\text{ins} = \int_0^1 f_\theta(\tilde{x}(t))\,dt \]
        \small im \textbf{większe} AUC, tym lepsza mapa istotności
      \end{block}
    \end{column}
  \end{columns}
  \vspace{8pt}
  \begin{alertblock}{Kluczowe założenie}
    Obie metryki zakładają, że \textbf{zastąpione piksele nie zawierają informacji}
    o klasie — co zależy od strategii zastępowania!
  \end{alertblock}
\end{frame}

\begin{frame}{Luka badawcza}
  \begin{columns}[T]
    \begin{column}{.55\textwidth}
      \textbf{Problem:} Usunięte piksele muszą być \emph{czymś} zastąpione.\\[6pt]
      Standardowe podejście: \textcolor{darkgray}{czarne piksele (zero)}\\[8pt]

      \begin{itemize}
        \item Tworzy \textbf{wejścia spoza rozkładu treningowego}\\
              (\textit{out-of-distribution, OoD})
        \item Model może tracić konfidencję z powodu \textbf{nienaturalności}
              obrazu, a nie braku cech dyskryminacyjnych
        \item Wyniki zależą od strategii — ale jak bardzo?
      \end{itemize}

      \vspace{6pt}
      \begin{exampleblock}{Future work (praca źródłowa)}
        \small ``\textit{One option would be to explore alternatives for replacing
        features in methods like Insertion/Deletion, such as leveraging a background
        dataset as in SHAP, which could provide new insight into feature attribution
        and interpretability.}''
      \end{exampleblock}
    \end{column}
    \begin{column}{.42\textwidth}
      \begin{tikzpicture}[scale=0.9]
        % Original
        \draw[fill=green!10, rounded corners=3pt] (0,4) rectangle (3.5,5);
        \node[font=\scriptsize] at (1.75,4.5) {Oryginalny obraz};
        % Arrow
        \draw[->, thick] (1.75,4) -- (1.75,3.4);
        % Replacement strategies
        \draw[fill=orange!20, rounded corners=3pt] (0,3.4) rectangle (1.5,2.5);
        \node[font=\tiny, align=center] at (0.75,2.95) {Czarne\\piksele};
        \draw[fill=blue!10, rounded corners=3pt] (2,3.4) rectangle (3.5,2.5);
        \node[font=\tiny, align=center] at (2.75,2.95) {Blur /\\Mean BG};
        \draw[->, thick] (0.75,2.5) -- (0.75,1.9);
        \draw[->, thick] (2.75,2.5) -- (2.75,1.9);
        % OoD
        \draw[fill=red!15, rounded corners=3pt] (0,1.9) rectangle (1.5,1.1);
        \node[font=\tiny, accentred, align=center] at (0.75,1.5) {Duże\\OoD};
        \draw[fill=green!15, rounded corners=3pt] (2,1.9) rectangle (3.5,1.1);
        \node[font=\tiny, successgreen, align=center] at (2.75,1.5) {Małe\\OoD};
        % Question
        \node[font=\scriptsize\bfseries, mainblue] at (1.75,0.6) {Która jest lepsza?};
      \end{tikzpicture}
    \end{column}
  \end{columns}
\end{frame}

% ============================================================
\section{Cel i pytania badawcze}
% ============================================================

\begin{frame}{Cel pracy i pytania badawcze}
  \begin{block}{Główny cel pracy}
    Empiryczna analiza wpływu strategii zastępowania pikseli na wyniki
    metryk Insertion/Deletion dla czterech metod XAI zastosowanych do
    klasyfikatora CNN.
  \end{block}

  \vspace{10pt}
  \begin{enumerate}
    \item[\textbf{P1}] Jak bardzo różnią się wartości AUC$_\text{ins}$ i
          AUC$_\text{del}$ w zależności od strategii zastępowania?
    \vspace{4pt}
    \item[\textbf{P2}] Czy strategia zastępowania zmienia \textbf{ranking} metod XAI?
          (korelacja rang Kendalla)
    \vspace{4pt}
    \item[\textbf{P3}] Która strategia jest \textbf{najmniej} podatna na problem OoD?
          (entropia predykcji na w pełni zastąpionych obrazach)
    \vspace{4pt}
    \item[\textbf{P4}] Czy generatywne uzupełnianie (inpainting)
          dostarcza bardziej \textbf{wiarygodnych} metryk niż proste baseline'y?
  \end{enumerate}
\end{frame}

% ============================================================
\section{Dane i model}
% ============================================================

\begin{frame}{Dane eksperymentalne i model}
  \begin{columns}[T]
    \begin{column}{.48\textwidth}
      \begin{block}{Dataset: ImageNette2-160}
        \begin{itemize}
          \item Podzbiór ImageNet — 10 klas
          \item \textasciitilde 13 000 obrazów treningowych
          \item \textasciitilde 500 obrazów walidacyjnych
          \item Rozdzielczość: $224\times224$ (resize)
          \item Klasy: tench, springer, cassette player,\\
                chain saw, church, French horn,\\
                garbage truck, gas pump,\\
                golf ball, parachute
        \end{itemize}
      \end{block}
      \vspace{4pt}
      \textbf{Do eksperymentów:} 50 obrazów (5 / klasę)\\
      \textbf{Zbiór tła (background):} 20 obrazów treningowych (2 / klasę)
    \end{column}
    \begin{column}{.48\textwidth}
      \begin{block}{Model: ResNet50}
        \begin{itemize}
          \item Pretrenowany na ImageNet (torchvision)
          \item Freeze — brak finetuningu
          \item Top-1 acc. na ImageNet: \textbf{76.1\%}
          \item Użyty do: forward pass + saliency
        \end{itemize}
      \end{block}
      \vspace{6pt}
      \begin{block}{Uzasadnienie wyboru}
        \begin{itemize}
          \item Popularny model referencyjny w literaturze XAI
          \item Dobrze zrozumiana architektura
          \item Łatwy dostęp do warstwy \texttt{layer4} (GradCAM)
        \end{itemize}
      \end{block}
    \end{column}
  \end{columns}
\end{frame}

% ============================================================
\section{Metody XAI}
% ============================================================

\begin{frame}{Metody XAI — przegląd}
  \begin{columns}[T]
    \begin{column}{.48\textwidth}
      \begin{block}{GradCAM \small(Selvaraju et al., 2017)}
        \begin{itemize}
          \item Gradient klasy względem ostatniej warstwy konwolucyjnej
          \item \textbf{Szybka}, działa lokalnie w sieci
          \item Niska rozdzielczość mapy — upsample do $224\times224$
        \end{itemize}
      \end{block}
      \vspace{6pt}
      \begin{block}{Integrated Gradients \small(Sundararajan et al., 2017)}
        \begin{itemize}
          \item Całka gradientu wzdłuż ścieżki od baseline do wejścia
          \item Spełnia aksjomaty atrybucji (completeness, sensitivity)
          \item Baseline: czarny obraz ($\mathbf{x}' = \mathbf{0}$)
        \end{itemize}
      \end{block}
    \end{column}
    \begin{column}{.48\textwidth}
      \begin{block}{LIME \small(Ribeiro et al., 2016)}
        \begin{itemize}
          \item Lokalne liniowe przybliżenie modelu
          \item Perturbacja superpikselami (SLIC segmentation)
          \item $n=500$ próbek, kolorowanie: czarne piksele
        \end{itemize}
      \end{block}
      \vspace{6pt}
      \begin{block}{SHAP — GradientExplainer \small(Lundberg, Lee, 2017)}
        \begin{itemize}
          \item Wartości Shapleya — sprawiedliwy podział atrybucji
          \item Wymaga \textbf{zbioru tła} (background dataset)
          \item 20 obrazów treningowych jako background
        \end{itemize}
      \end{block}
    \end{column}
  \end{columns}
  \vspace{6pt}
  \small
  Wszystkie metody zwracają mapę $\mathbf{S} \in [0,1]^{H \times W}$
  — znormalizowaną do przedziału $[0,1]$ przed dalszą analizą.
\end{frame}

% ============================================================
\section{Strategie zastępowania pikseli}
% ============================================================

\begin{frame}{Strategie zastępowania pikseli — przegląd}
  \renewcommand{\arraystretch}{1.4}
  \begin{tabular}{>{\bfseries\small}p{2.6cm}
                  >{\small}p{4.5cm}
                  >{\small}p{3.8cm}}
    \toprule
    \textbf{Strategia} & \textbf{Opis} & \textbf{Zalety / wady} \\
    \midrule
    \rowcolor{lightgray}
    Black (zero) &
      Piksele $\leftarrow 0$ (w przestrzeni znormalizowanej) &
      Prosta, szybka \textcolor{accentred}{/ duże OoD} \\

    Gaussian Blur &
      Piksele $\leftarrow$ wygładzony oryginał ($\sigma=10$) &
      Zachowuje strukturę \textcolor{accentred}{/ może zachować cechy} \\

    \rowcolor{lightgray}
    Mean Background &
      Piksele $\leftarrow$ średni piksel ze zbioru treningowego &
      Zbliżone do SHAP \textcolor{successgreen}{/ in-distribution} \\

    Random Noise &
      Piksele $\leftarrow \mathcal{N}(0, \sigma^2)$ &
      Stochastyczna \textcolor{accentred}{/ najmniej naturalna} \\

    \rowcolor{lightgray}
    CV2 Inpainting &
      Algorytm Telea — uzupełnia z sąsiadów &
      Spójna wizualnie \textcolor{successgreen}{/ szybka} \\

    SD Inpainting &
      Stable Diffusion 2 (opcjonalne) &
      Fotorealistyczna \textcolor{accentred}{/ wolna, stochastyczna} \\
    \bottomrule
  \end{tabular}

  \vspace{6pt}
  \begin{alertblock}{Hipoteza robocza}
    Strategie bliższe rozkładowi treningowemu (Mean BG, CV2 Inpainting)
    dadzą \textbf{bardziej stabilne i wiarygodne} metryki niż Black / Noise.
  \end{alertblock}
\end{frame}

\begin{frame}{Algorytm CV2 Inpainting (Telea, 2004)}
  \begin{columns}[T]
    \begin{column}{.5\textwidth}
      \begin{itemize}
        \item Uzupełnia maskowane obszary propagując z granic
        \item Oparty na równaniach Naviera-Stokesa
        \item Parametr: promień inpaintingu ($r=3$ px)
        \item Dylatacja maski (kernel $3\times3$, 1 iteracja)
              zapobiega artefaktom na krawędzi
      \end{itemize}
      \vspace{6pt}
      Złożoność: $\mathcal{O}(|\Omega| \cdot r^2)$ gdzie $\Omega$ = obszar maski
    \end{column}
    \begin{column}{.47\textwidth}
      \begin{block}{Algorytm (schemat)}
        \small
        \begin{enumerate}
          \item Wyznacz maskę $M$ z mapy istotności
          \item Dylatacja: $M' = M \oplus B_{3\times3}$
          \item Inpainting: $I' = \text{Telea}(I, M')$
          \item Sklejanie: $I_\text{out}[M=0] = I'[M=0]$
        \end{enumerate}
      \end{block}
    \end{column}
  \end{columns}
\end{frame}

% ============================================================
\section{Projekt eksperymentów}
% ============================================================

\begin{frame}{Projekt eksperymentów}
  \begin{block}{Schemat eksperymentu}
    $N_\text{img} = 50$ obrazów $\times$
    $N_\text{XAI} = 4$ metody $\times$
    $N_\text{strat} = 5$ strategii
    $= \mathbf{1000}$ obserwacji (AUC$_\text{ins}$ + AUC$_\text{del}$)
  \end{block}

  \vspace{8pt}
  \begin{columns}[T]
    \begin{column}{.5\textwidth}
      \textbf{Dla każdego obrazu $i$, metody $m$, strategii $s$:}
      \begin{enumerate}
        \item Oblicz mapę istotności $\mathbf{S}_{i,m}$
        \item Posortuj piksele malejąco: $\pi = \text{argsort}(\mathbf{S})$
        \item Dla $t \in \{0, \frac{1}{T}, \ldots, 1\}$ ($T=20$ kroków):
        \begin{itemize}
          \item Insertion: zastąp wszystko oprócz top-$t$
          \item Deletion: zastąp top-$t$
          \item Zmierz $f_\theta(x_t)[y_i]$
        \end{itemize}
        \item Oblicz AUC metodą trapezów
      \end{enumerate}
    \end{column}
    \begin{column}{.47\textwidth}
      \begin{tikzpicture}[scale=0.85]
        \begin{axis}[
          xlabel={\small Ułamek pikseli},
          ylabel={\small Konfidencja},
          xmin=0, xmax=1, ymin=0, ymax=1,
          width=5.5cm, height=4cm,
          legend style={font=\tiny, at={(0.98,0.02)}, anchor=south east},
          grid=major, grid style={dotted, gray!40},
          tick label style={font=\tiny}
        ]
        % Insertion (ideal)
        \addplot[color=successgreen, thick, smooth]
          coordinates {(0,0.05)(0.1,0.2)(0.2,0.45)(0.3,0.6)
                       (0.5,0.75)(0.7,0.85)(1,0.88)};
        \addlegendentry{Insertion (↑ better)}
        % Deletion (ideal)
        \addplot[color=accentred, thick, smooth]
          coordinates {(0,0.88)(0.1,0.72)(0.2,0.55)(0.3,0.4)
                       (0.5,0.25)(0.7,0.12)(1,0.05)};
        \addlegendentry{Deletion (↓ better)}
        % AUC fill
        \addplot[fill=successgreen, fill opacity=0.1, draw=none]
          coordinates {(0,0)(0.1,0.2)(0.2,0.45)(0.3,0.6)
                       (0.5,0.75)(0.7,0.85)(1,0.88)(1,0)(0,0)};
        \end{axis}
      \end{tikzpicture}
    \end{column}
  \end{columns}
\end{frame}

% ============================================================
\section{Kryteria porównawcze}
% ============================================================

\begin{frame}{Kryteria porównawcze}
  \renewcommand{\arraystretch}{1.35}
  \small
  \begin{tabular}{>{\bfseries}p{3.6cm}
                  >{\raggedright\arraybackslash}p{4.5cm}
                  >{\raggedright\arraybackslash}p{3.4cm}}
    \toprule
    \textbf{Kryterium} & \textbf{Miara} & \textbf{Interpretacja} \\
    \midrule
    \rowcolor{lightgray}
    Wierność – Insertion &
      $\text{AUC}_\text{ins} \in [0,1]$ &
      Wyższe = lepsza mapa XAI \\

    Wierność – Deletion &
      $\text{AUC}_\text{del} \in [0,1]$ &
      Niższe = lepsza mapa XAI \\

    \rowcolor{lightgray}
    Wrażliwość na OoD &
      Konfidencja i entropia na w pełni zastąpionych obrazach &
      Niższa konfidencja = mniej OoD \\

    Stabilność rankingu &
      Kendall's $\tau$ między parami strategii &
      $\tau \approx 1$ = strategie są zgodne \\

    \rowcolor{lightgray}
    Jakość wizualna &
      SSIM między oryginałem a zastąpionym (przed maskowaniem) &
      Wyższy SSIM = mniej artefaktów \\

    Czas obliczeń &
      Czas na 1 obraz $\times$ 1 metodę XAI [s] &
      Praktyczna wykonalność \\
    \bottomrule
  \end{tabular}

  \vspace{8pt}
  \begin{block}{Złożony wskaźnik jakości (propozycja)}
    \[ Q = \text{AUC}_\text{ins} - \text{AUC}_\text{del} + \alpha \cdot (1 - P_\text{OoD}) \]
    gdzie $\alpha=0.5$ balansuje wierność i OoD-awareness.
  \end{block}
\end{frame}

% ============================================================
\section{Oczekiwane wyniki}
% ============================================================

\begin{frame}{Oczekiwane wyniki i hipotezy}
  \begin{columns}[T]
    \begin{column}{.5\textwidth}
      \begin{block}{H1 — Różnica AUC między strategiami}
        Spodziewamy się \textbf{istotnych statystycznie} różnic AUC
        między Black a CV2 Inpainting ($>0.05$ punktu AUC)
      \end{block}
      \vspace{6pt}
      \begin{block}{H2 — Zmiana rankingu XAI}
        Strategia zastępowania zmienia ranking metod XAI
        (Kendall's $\tau < 0.8$ między Black a CV2 Inpainting)
      \end{block}
    \end{column}
    \begin{column}{.47\textwidth}
      \begin{block}{H3 — Inpainting redukuje OoD}
        CV2 Inpainting daje niższą konfidencję na w pełni
        zastąpionych obrazach niż Black i Noise
      \end{block}
      \vspace{6pt}
      \begin{block}{H4 — LIME najbardziej wrażliwy}
        LIME (oparty na superpikselach) będzie wykazywał
        najwyższą wariancję AUC między strategiami,
        bo jego segmentacja zakłada konkretny sposób perturbacji
      \end{block}
    \end{column}
  \end{columns}
\end{frame}

% ============================================================
\section{Harmonogram i podsumowanie}
% ============================================================

\begin{frame}{Harmonogram prac}
  \renewcommand{\arraystretch}{1.3}
  \begin{tabular}{>{\bfseries}p{2.8cm}
                  >{\raggedright\arraybackslash}p{8.5cm}}
    \toprule
    \textbf{Faza} & \textbf{Zadania} \\
    \midrule
    Grudzień 2026 &
      Implementacja pipeline'u eksperymentalnego (notebook)\\
      & Testy i walidacja implementacji na 10 obrazach \\

    Styczeń 2027 &
      Pełne eksperymenty (50 obrazów × 4 metody × 5 strategii)\\
      & Analiza wyników, wykresy, tabela podsumowująca \\

    Luty 2027 &
      Eksperymenty z SD inpainting (opcjonalnie) \\
      & Pisanie rozdziału metodologicznego pracy \\

    Marzec 2027 &
      Pisanie rozdziału z wynikami i wnioskami \\
      & Redakcja końcowa, oddanie pracy \\
    \bottomrule
  \end{tabular}
\end{frame}

\begin{frame}{Podsumowanie}
  \begin{columns}[T]
    \begin{column}{.55\textwidth}
      \textbf{Wkład pracy (contribution):}
      \begin{itemize}
        \item Pierwsza systematyczna analiza porównawcza
              5 strategii zastępowania pikseli w metrykach Ins./Del.
        \item Empiryczne potwierdzenie/obalenie hipotezy OoD
        \item Rekomendacja strategii dla przyszłych badań XAI
        \item Publiczny notebook Google Colab z reprodukowalnym kodem
      \end{itemize}
    \end{column}
    \begin{column}{.42\textwidth}
      \begin{block}{Kluczowe metryki eksperymentu}
        \begin{itemize}
          \item \textbf{50} obrazów, 10 klas
          \item \textbf{4} metody XAI
          \item \textbf{5} strategii zastępowania
          \item \textbf{20} kroków per krzywa
          \item $\approx$ \textbf{2000} pomiarów AUC
        \end{itemize}
      \end{block}
    \end{column}
  \end{columns}

  \vspace{10pt}
  \begin{center}
    \textbf{\Large Dziękuję za uwagę!}\\[6pt]
    \small Pytania? Uwagi?
  \end{center}
\end{frame}

% ============================================================
% Slajdy uzupełniające (backup)
% ============================================================

\appendix

\begin{frame}[allowframebreaks]{Literatura}
  \small
  \begin{thebibliography}{99}

    \bibitem{petsiuk2018}
    Petsiuk, V., Das, A., \& Saenko, K. (2018).
    \textit{RISE: Randomized Input Sampling for Explanation of Black-box Models}.
    BMVC.

    \bibitem{selvaraju2017}
    Selvaraju, R.R., et al. (2017).
    \textit{Grad-CAM: Visual Explanations from Deep Networks via Gradient-based Localization}.
    ICCV.

    \bibitem{sundararajan2017}
    Sundararajan, M., Taly, A., \& Yan, Q. (2017).
    \textit{Axiomatic Attribution for Deep Networks}.
    ICML.

    \bibitem{ribeiro2016}
    Ribeiro, M.T., Singh, S., \& Guestrin, C. (2016).
    \textit{"Why Should I Trust You?": Explaining the Predictions of Any Classifier}.
    KDD.

    \bibitem{lundberg2017}
    Lundberg, S.M., \& Lee, S.I. (2017).
    \textit{A Unified Approach to Interpreting Model Predictions}.
    NeurIPS.

    \bibitem{agarwal2020}
    Agarwal, C., \& Nguyen, A. (2020).
    \textit{Explaining Image Classifiers by Removing Input Features Using Generative Models}.
    ACCV.

    \bibitem{back2baseline2025}
    Back to the Baseline: Examining Baseline Effects on Explainability Metrics.
    arXiv:2512.11433.

    \bibitem{decoupling2024}
    Blücher, S. et al. (2024).
    \textit{Decoupling Pixel Flipping and Occlusion Strategy for Consistent XAI Benchmarks}.
    arXiv:2401.06654.

    \bibitem{lili2025}
    Inpainting Insights: Elevating Visual XAI with Photorealistic Perturbations.
    arXiv:2607.15482.

    \bibitem{telea2004}
    Telea, A. (2004).
    \textit{An Image Inpainting Technique Based on the Fast Marching Method}.
    Journal of Graphics Tools, 9(1), 23–34.

  \end{thebibliography}
\end{frame}

\begin{frame}{Uzupełnienie: Stable Diffusion Inpainting}
  \begin{columns}[T]
    \begin{column}{.5\textwidth}
      \textbf{Model:} \texttt{stabilityai/stable-diffusion-2-inpainting}
      \begin{itemize}
        \item U-Net z warunkowym wejściem maski
        \item 20 kroków DDIM (kompromis prędkość/jakość)
        \item Dylatacja maski (8px) — kluczowa dla usunięcia artefaktów
        \item Prompt: ``a natural scene'' (neutralny)
      \end{itemize}
      \vspace{6pt}
      \textbf{Ograniczenia:}
      \begin{itemize}
        \item Stochastyczność — każde wywołanie daje inny wynik
        \item Semantic leakage — może halucynować cechy klasy
        \item Wymaga $>10$ GB VRAM (Colab A100)
        \item $\approx 3$ s / obraz (vs $<0.01$ s dla CV2)
      \end{itemize}
    \end{column}
    \begin{column}{.47\textwidth}
      \begin{block}{Procedura}
        \begin{enumerate}
          \item Mapa istotności $\rightarrow$ maska binarna
          \item Dylatacja maski (kernel 8px)
          \item Przekaż do SD2: obraz + maska + prompt
          \item Próbkuj przez 20 kroków (DDIM)
          \item Sklejanie: zamień tylko zamaskowane piksele
        \end{enumerate}
      \end{block}
      \vspace{4pt}
      \small
      Semantic leakage to kluczowe wyzwanie: SD może rekonstruować
      obiekt na podstawie kontekstu, fałszywie zawyżając konfidencję.
    \end{column}
  \end{columns}
\end{frame}

\end{document}
